\section{Conclusion}\label{Conclusion}
Reliable state estimation is a fundamental requirement for safe and effective autonomous mobile manipulation. 
This study investigated the comparative performance of three localization algorithms: Dead Reckoning (DR), the Extended Kalman Filter (EKF), and the Particle Filter (PF). 
By leveraging ROS2 and Gazebo, the algorithms were deployed and evaluated on a simulated differential drive mobile manipulator within a realistic factory environment. 
The algorithms utilized a unicycle kinematic model with injected odometry noise, incorporating 2-D LiDAR measurements through Iterative Closest Point (ICP) scan matching for the EKF and a Likelihood Field model for the PF. 
The primary objective was to evaluate the trade-offs between local tracking accuracy, computational efficiency, and global recovery capabilities.

The results supported the initial hypotheses regarding the strengths and weaknesses of Gaussian and non-parametric filters. 
Under baseline navigation conditions, the EKF demonstrated high tracking accuracy while maintaining a low computational footprint. 
This makes the EKF highly suitable for resource-constrained systems provided the initial pose is well defined and tracking remains continuous. 
However, both the EKF and DR approaches completely failed to recover from catastrophic tracking failures during the kidnapped robot scenarios. 
Conversely, the PF successfully resolved global displacements by maintaining multiple pose hypotheses until it reconverged on the true state. 
Increasing the number of particles in the PF predictably improved both baseline accuracy and kidnap recovery times, but this came at the cost of a significantly higher computational overhead.

While these findings provide valuable insights into 2-D planar localization, the study possesses several limitations. 
The experiments were conducted in a simulated environment, which may not perfectly capture the complex, non-Gaussian sensor anomalies present in physical LiDAR and wheel encoders. 
Additionally, the factory environment was entirely static. 
The presence of dynamic obstacles, such as human workers or other robots, would introduce unmodeled occlusions that could degrade the performance of the ICP and Likelihood Field measurement updates. 

Future work should be focused on comparing how these algorithms perform as an input into local motion planning algorithms.
For instance, although the PF has a higher accuracy it can be known to jitter.
Does this jitter cause issues with the position controller and local motion planner that are downstream?
Additionally, applying these algorithms to a sloped map will require fundamental changes to the underlying system model. 
The state vector must be expanded from three degrees of freedom (planar coordinates and yaw) to a full six degrees of freedom to account for pitch, roll, and elevation changes. 
Consequently, the 2-D unicycle kinematic models must be replaced with 3-D motion models that account for gravitational slip and complex terrain geometry. 
Finally, the sensor model must transition from planar 2-D LiDAR projections and 2-D OccupancyGrids to full 3-D point cloud registration techniques using volumetric maps, heavily increasing the computational requirements for both the EKF and the PF.